\section{Architettura e scelte progettuali}
\label{sec:architettura}
\subsection{Flusso dei dati e responsabilità}
La Fig.~\ref{fig:architecture} rappresenta la pipeline. Ciascun sito comprende
un broker Mosquitto e un Edge Gateway; il Simulator corrispondente pubblica
JSON sul topic MQTT \code{sensors/<sensorID>/telemetry}. Il gateway valida le
misure e genera un \code{EdgeAggregate} per zona e finestra non vuota di cinque
minuti. Il record contiene eventi, conteggi validi e non validi, somma, media,
minimo e massimo per ciascuna grandezza.

\begin{figure}[t]
\centering
\begin{tikzpicture}[
  node distance=4mm,
  box/.style={draw,rounded corners=1pt,align=center,text width=6.8cm,
    minimum height=7mm,font=\footnotesize,inner sep=3pt},
  flow/.style={-{Latex[length=1.8mm]},semithick}]
\node[box,fill=gray!7] (sim) {13 Simulator: replay di 150 sensori};
\node[box,below=of sim,fill=blue!6] (edge) {13 siti Edge: Mosquitto + Gateway\\
  aggregazione locale 5 min e watermark per Edge};
\node[box,below=of edge] (k1) {Kafka: \texttt{edge-aggregates}\\6 partizioni, chiave \texttt{edge\_id}};
\node[box,below=of k1,fill=orange!9] (workers) {$W$ Cloud Worker, stesso consumer group\\
  finestre 15 min, minimo dei watermark per partizione};
\node[box,below=of workers] (k2) {Kafka: \texttt{cloud-partition-aggregates}\\
  1 partizione fisica, 6 sorgenti logiche};
\node[box,below=of k2,fill=orange!9] (global) {Global Aggregator\\
  riduzione degli stessi intervalli di 15 min};
\node[box,below=of global,fill=gray!7] (sink) {Sink: log JSON oppure PostgreSQL};
\foreach \a/\b in {sim/edge,edge/k1,k1/workers,workers/k2,k2/global,global/sink}
  \draw[flow] (\a) -- (\b);
\end{tikzpicture}
\caption{Architettura logica. I due topic risiedono sullo stesso broker Kafka.
Le repliche dei worker non modificano le sei sorgenti logiche del Global.}
\label{fig:architecture}
\end{figure}

Kafka disaccoppia la produzione Edge dal consumo Cloud e mantiene i record in
un log persistente~\cite{kafka}. Il topic \code{edge-aggregates} ha sei
partizioni. La chiave \code{edge\_id} e il partizionatore hash mantengono
\code{EdgeAggregate}, \code{EdgeWatermark} e marker terminale dello stesso
Edge nella medesima partizione e nello stesso ordine. Il numero di partizioni
e la membership degli Edge rimangono invariati per l'intera run.

I Cloud Worker condividono un consumer group: ciascuna partizione è assegnata
a un solo membro del gruppo. La topologia Edge è proiettata sulle partizioni
con lo stesso hash del producer; ogni worker mantiene uno stato separato per
ciascuna partizione assegnata, indipendente dalla propria identità. I contributi
Edge sono combinati su finestre di quindici minuti. Gli output sono
\code{CloudPartitionAggregate}, identificati da partizione sorgente e confini
temporali, pubblicati sul secondo topic Kafka. Il Global Aggregator combina i
partial con la stessa coppia di confini e produce il risultato nel sink.

\subsection{Aggregazione componibile}
Per una metrica $m$, ogni aggregato conserva il numero di misure valide $n_m$,
il numero di misure non valide $u_m$ e la somma $s_m$, oltre agli estremi.
La combinazione di contributi disgiunti $j$ applica
\begin{equation}
n_m=\sum_j n_{m,j},\qquad u_m=\sum_j u_{m,j},\qquad
s_m=\sum_j s_{m,j}.
\end{equation}
La media si ottiene da $s_m/n_m$ se $n_m>0$; altrimenti media ed estremi sono
nulli. Minimo e massimo si combinano considerando soltanto contributi validi.
La somma dei conteggi validi e non validi deve coincidere con il numero di
eventi dell'aggregato per ciascuna metrica.

Conservare la somma evita la media non pesata delle medie locali, che sarebbe
errata in presenza di numerosità differenti. L'operazione è componibile in
aritmetica esatta; con valori floating point l'equivalenza numerica va
verificata con una tolleranza. I cinque minuti Edge dividono esattamente i
quindici minuti Cloud, per cui ogni aggregato locale appartiene a una sola
finestra Cloud. Un aggregato che attraversa un confine Cloud viene rifiutato.

\subsection{Progresso, contributi vuoti e completamento}
L'Edge mantiene una finestra corrente in event time. L'arrivo di un evento
della finestra successiva determina la pubblicazione della precedente, seguita
da \code{EdgeWatermark(e,t)}. Tale controllo certifica che l'Edge $e$ non
pubblicherà successivamente aggregati con \code{window\_end} minore o uguale
a $t$. Il watermark è monotono. L'EOS ricevuto dal Simulator chiude l'ultima
finestra; dopo averne completato la scrittura Kafka, l'Edge pubblica il distinto
marker \code{EdgeEndOfInput}. Gli eventi appartenenti a finestre già chiuse
sono scartati e conteggiati; il contratto richiede shard ordinati temporalmente.

Una partizione Cloud può contenere Edge con velocità differenti. Il worker
memorizza, per ogni Edge atteso nella partizione $P$, il watermark
$W_e$ e il flag terminale. Finché un Edge attivo non ha prodotto un watermark,
non è possibile certificarne il progresso. La frontiera locale è
\begin{equation}
W_P=\min_{e\in E_P,\ e\ \mathrm{attivo}} W_e .
\end{equation}
Solo le finestre con fine non successiva a $W_P$ sono definitive. Un Edge
terminato non partecipa più al minimo; la cessazione temporanea del traffico,
invece, non equivale alla conclusione della sorgente.

Il worker pubblica prima i partial definitivi e poi il controllo di progresso
\code{PartitionProgress(P,t)}, con $t=W_P$. Quando tutti gli Edge di $P$ terminano,
pubblica in ordine gli ultimi partial e infine \code{PartitionEndOfReplay(P)};
una partizione senza Edge produce direttamente tale controllo. Un aggregato
ricevuto dietro un watermark già certificato viola il protocollo e invalida la
run: una finestra emessa non viene riaperta. L'ordine fra dati e controlli
permette una rappresentazione sparsa: l'assenza di un partial, dopo il relativo
progresso o EOS, equivale a un contributo zero.

È stata considerata anche un'alternativa in cui il Cloud conoscesse soltanto
le partizioni Kafka, senza la membership degli Edge, e stimasse la frontiera
come massimo event time osservato meno un ritardo. La soluzione ridurrebbe
l'accoppiamento con la topologia, ma un Edge veloce potrebbe anticipare la
chiusura rispetto a uno lento; inoltre il ritardo sarebbe un parametro euristico.
Si è quindi preferita la certificazione esplicita degli Edge e il minimo per
partizione, privilegiando la correttezza nel deployment statico adottato.

Il Global conosce soltanto le sei partizioni sorgente. Per una finestra attende
da ciascuna un partial oppure un controllo di progresso/EOS che certifichi il
contributo zero. Non introduce nuove finestre, non ricalcola watermark e non
usa timeout per dichiarare completi risultati incompleti. Il replay termina
dopo gli EOS di tutte le partizioni e l'emissione delle finestre pendenti.

\subsection{Parallelismo indipendente dal risultato}
Con $W=1,2,4,6$ cambia il numero di esecutori, non quello delle partizioni
sorgente (Tab.~\ref{tab:workers}). Anche con un solo worker la pipeline passa
dal secondo topic e dal Global. In presenza di dati in tutte le partizioni si
producono sei partial per finestra; in caso contrario alcuni contributi sono
certificati come zero. Il Global non aggrega per identità del worker.

\begin{table}[t]
\caption{Distribuzione bilanciata delle sei partizioni. L'assegnazione
effettiva degli identificativi è decisa dal consumer group.}
\label{tab:workers}
\centering
\begin{tabular}{ccc}
\toprule
Worker & Partizioni per worker & Sorgenti logiche Global\\
\midrule
1 & 6 & 6\\
2 & 3, 3 & 6\\
4 & 2, 2, 1, 1 & 6\\
6 & 1, 1, 1, 1, 1, 1 & 6\\
\bottomrule
\end{tabular}
\end{table}

Ogni worker usa letture concorrenti ma un solo ciclo di elaborazione; il
parallelismo computazionale cresce quindi tramite processi aggiuntivi. Sei
partizioni costituiscono il limite del parallelismo utile del gruppo. La
distribuzione delle partizioni non garantisce un uguale costo: il numero di
Edge e il volume dei loro aggregati possono differire. Il broker condiviso e
il Global restano inoltre possibili colli di bottiglia.
